\documentclass{article}
\usepackage[UTF8]{ctex}
\usepackage{amsmath,amssymb,amsfonts}
\usepackage{enumitem}
\usepackage{geometry}
\geometry{a4paper, margin=1in}
% === 错题本排版优化（不参与题目标识） ===
\setlength{\parindent}{0pt}
\setlength{\parskip}{0.35em plus 0.10em minus 0.08em}
\setlength{\emergencystretch}{3em}
\setlist{itemsep=0.25em, topsep=0.40em, parsep=0pt}
\allowdisplaybreaks[2]
% === 错题本排版优化结束 ===
\begin{document}
\section*{华东师范大学2025年数学分析真题}
\section*{37．华东师范大学2025年数学分析试卷（回忆版）}

\subsection{题目1}
1． \(\lim _{n \rightarrow \infty} \sqrt[n]{\left(1+\frac{1}{2}\right)\left(1+\frac{2}{3}\right) \cdots\left(1+\frac{n-1}{n}\right)}=\) \(\_\_\_\_\)．

\subsection{题目2}
2． \(\int_{0}^{\pi} \frac{x \sin ^{2} x}{1+\cos ^{2} x} \mathrm{~d} x=\) \(\_\_\_\_\)．

\subsection{题目3}
3． \(\lim _{x \rightarrow \infty}\left(\cos \frac{a}{x}\right)^{x^{2}}=\) \(\_\_\_\_\)．

\subsection{题目4}
4．\(\sum_{n=1}^{\infty}(-1)^{n-1} \frac{\sin n x}{n}\) 是 \(\_\_\_\_\)在 \((-\pi, \pi)\) 上的傅里叶级数．

\subsection{题目5}
5．求和函数 \(\sum_{n=0}^{\infty} \frac{n^{2}+1}{n!2^{n}} x^{n}=\) \(\_\_\_\_\)．

\subsection{题目6}
6．由柱面坐标表示的曲面 \(z=z(r, \theta)\) ，其上任意一条光滑曲线的弧长微分 \(\mathrm{d} s\) 满足 \(\mathrm{d} s^{2}=\)\\
\(\_\_\_\_\) \(\mathrm{d} r^{2}+\) \(\_\_\_\_\) \(\mathrm{d} r \mathrm{~d} \theta+\) \(\_\_\_\_\) \(\mathrm{d} \theta^{2}\) ．

\subsection{题目7}
7．求第二型曲面积分 \(\iint_{S} \frac{a^{2} x \mathrm{~d} y \mathrm{~d} z+(z+a)^{2} \mathrm{~d} x \mathrm{~d} y}{\left(x^{2}+y^{2}+z^{2}\right)^{\frac{3}{2}}}=\) \(\_\_\_\_\)，其中 \(S\) 为 \(z=-\sqrt{a^{2}-x^{2}-y^{2}}\) ，取上侧。

\subsection{题目8}
1．解答如下问题：\\
（1）\(f(x)\) 为定义在 \((a, b)\) 上的函数，若对任意的 \(x_{0} \in(a, b)\) ，有 \(\lim _{x \rightarrow x_{0}} f(x)\) 存在，任取 \(\left\{a_{n}\right\} \subset(a, b)\) 为Cauchy 列，那么是否有 \(\left\{f\left(a_{n}\right)\right\}\) 为 Cauchy 列？\\
（2）\(f(x)\) 为定义在 \((a, b)\) 上的函数，若对 \((a, b)\) 中任意的 Cauchy 列 \(\left\{a_{n}\right\},\left\{f\left(a_{n}\right)\right\}\) 也是 Cauchy 列，证明：\(f(x)\) 在 \((a, b)\) 上一致连续．

\subsection{题目9}
2．设 \(f(x)\) 在 \([a, b]\) 上可微．

\begin{itemize}
  \item[（i）] \(f^{\prime}(x)\) 在 \([a, b]\) 上连续．
  \item[（ii）] 对任意的 \(\varepsilon>0\) ，存在 \(\delta>0\) ，使得当 \(0<|h|<\delta\) 时，有 \(\left|\frac{f(x+h)-f(x)}{h}-f^{\prime}(x)\right|<\varepsilon\) 对任意的 \(x \in[a, b]\)成立。
\end{itemize}

讨论（i）能否推（ii）？（ii）能否推（i）？

\subsection{题目10}
3．设 \(\left\{f_{n}(x)\right\}\) 为 \([a, b]\) 上连续的函数列，且 \(\sum_{n=1}^{\infty} f_{n}(x)\) 在 \((a, b)\) 上一致收敛。

\begin{itemize}
  \item[（1）] 证明：\(\sum_{n=1}^{\infty} f_{n}(a), \sum_{n=1}^{\infty} f_{n}(b)\) 收敛。
  \item[（2）] 证明：\(\sum_{n=1}^{\infty} f_{n}(x)\) 在 \([a, b]\) 上一致收敛。
\end{itemize}

\subsection{题目11}
4．设 \(f(x)\) 为 \((-\infty,+\infty)\) 上的函数，且在任意有限区间上可积，若 \(\int_{-\infty}^{+\infty}|f(x)| \mathrm{d} x\) 收敛，证明：\(F(x)=\) \(\int_{-\infty}^{+\infty} f(t) \cos 2 \pi t x \mathrm{~d} t\) 在 \((-\infty,+\infty)\) 上一致连续．

\subsection{题目12}
5．解答如下问题：

\begin{itemize}
  \item[（1）] 对于 \(\sum_{n=1}^{\infty} a_{n}\) ，若对任意子列 \(\left\{a_{n_{k}}\right\}\) ，均有 \(\sum_{k=1}^{\infty} a_{n_{k}}\) 收敛，那么 \(\sum_{n=1}^{\infty} a_{n}\) 是否绝对收敛？
  \item[（2）] 对于任意正整数 \(i, j\) ，均有 \(\sum_{n=1}^{\infty} a_{i+n j}\) 收敛，是否有 \(\sum_{n=1}^{\infty} a_{n}\) 绝对收敛？
\end{itemize}

\subsection{题目13}
6．设 \(V, U\) 分别是 \(\mathbb{R}^{m}, \mathbb{R}^{n}\) 上的开集，定义向量函数

\[
f: V \mapsto U,\left(x_{1}, \cdots, x_{m}\right) \rightarrow\left(f_{1}\left(x_{1}, \cdots, x_{m}\right), \cdots, f_{n}\left(x_{1}, \cdots, x_{m}\right)\right) .
\]

称 \(f\) 在点 \(x_{0}\) 可微，如果存在 \(m \times n\) 矩阵 \(A\) ，满足

\[
f\left(x_{0}+h\right)=f\left(x_{0}\right)+h A+o(\|h\|) .
\]

其中 \(\mathbf{h}=\left(h_{1}, \cdots, h_{m}\right), o(\|\mathbf{h}\|)\) 为各元素均为 \(\|\mathbf{h}\|\) 高阶无穷小的 \(m\) 维行向量。此外，定义 Jacobi 阵为 \(D_{x_{0}}^{f}=\left.\left(\frac{\partial f_{j}}{\partial x_{i}}\right)_{m \times n}\right|_{x=x_{0}}\) 。

\begin{itemize}
  \item[（1）] 证明：若 \(f\) 在 \(x_{0}\) 处可微，则 \(D_{x_{0}}^{f}=A\) 。
  \item[（2）] 进一步，若 \(g\) 在 \(f\left(x_{0}\right)\) 处可微，则 \(g \circ f\) 在 \(x_{0}\) 处可微，且有 \(A=D_{f\left(x_{0}\right)}^{g} D_{x_{0}}^{f}\) ．
  \item[（3）] 未知…
\end{itemize}
\end{document}
